\begin{abstract}
Electronic Health Records (EHRs) contain a mix of structured tables and unstructured notes that are difficult to query using keyword search or SQL alone. This dissertation presents the design and evaluation of a Retrieval-Augmented Generation (RAG) chatbot built on the MIMIC-IV dataset. The system integrates a locally hosted Large Language Model (LLM) with FAISS-based semantic retrieval to generate grounded, evidence-backed clinical answers. 

The pipeline covers data preprocessing, chunking, embedding construction, and retrieval, with multiple compact LLMs (Phi-3, Qwen, Gemma, Llama, DeepSeek) evaluated via Ollama. An API, command-line interface, and web frontend were developed to support modular access. A gold-standard set of clinical questions was used to measure performance across precision, recall, F1 score, retrieval latency, and hallucination rates. 

Results show that the RAG pipeline consistently outperformed a standalone LLM baseline, achieving higher factual accuracy and reducing unsupported claims. Trade-offs between embedding choice, model size, and response time were analysed, with optimal configurations identified for both accuracy-focused and efficiency-focused deployment scenarios. 

This work contributes a reproducible, locally deployable clinical RAG framework that demonstrates the feasibility of safe and privacy-preserving clinical decision support systems. Limitations include the lack of live clinical validation and reliance on synthetic or de-identified datasets. Future work should extend to multimodal data, stronger bias mitigation, and integration into real EHR workflows.
\end{abstract}
